\section{Safety and Provenance}
\sectionkicker{SAFETY}
\begin{wideflow}
{\Large\bfseries 安全性と来歴の5領域――来歴評価、公開ガードレール、透かし、指針、認可仕様}\par
\smallskip
{\color{SurveyMuted}5つの独立した領域を別々に読む。単一の安全性スコアには束ねない。}\par
\end{wideflow}

\subsection*{来歴評価――ProvenanceGuard}
ProvenanceGuardは、今週のチームブログ記事に載った出来事と、対象期間より前の4著者の論文（v1と7月26日のv2）だけで読む\cite{provguard}。v3の有無については何も言わない。ルーティングと自然言語推論、帰属づけと許可・遮断、RARRという構成で、281件の追跡のうち138件と139件を検出し、F1スコアで0.802と0.858、複数情報源の組み合わせで0.846、0.503、0.229と置かれる。ベースラインとの比較は支えに回し、修復やフォールバック、半々のプローブも添える。これは帰属関係を評価する層の話であり、実行規制や電子透かしの話と混ぜない。

\subsection*{来歴評価の測定条件}
測定値の読み方に注意が要る。281件は合成・匿名化の凍結済み医療追跡（個人情報なし）、266件・2325ラベルの部分集合、40件361主張のホールドアウト（LLM 2審判と優先づけ、人手専門家は361件のみ）という別の実験設定の数値であり、一つの優劣に畳んではならない。ブロックF1は0.802（適合率0.673・再現率0.993）、情報源正答率0.858/260、関係つき0.681、複数情報源では0.846・0.503・0.229と置かれる。対象は帰属の正しさに限り、臨床的安全性やパラメトリックな主張ではない。図表はテキストの仮置きとして読み、5.4節以降や付録の範囲を超えた部分は未読である。コードやポスター、リポジトリの変更提案は開いていない。来歴評価という機能に限った読みとする。

\subsection*{公開ガードレール――Open Agent Safety}
NVIDIAのOpen Agent Safety Platformは、公開されたエージェント見守りのリファレンスアーキテクチャである\cite{agentsafety}。OpenShellのOSSランタイムは隔離実行・方針実施・検証可能追跡を担い広く提供され、SentryはBlueField-4上の帯域外番犬としてミリ秒単位の隔離とDOCAによる計測・証明つき遠隔測定を添える参照設計である。Vera向けに最適化し、ArmやIntelにも拡張可能とされる。参照構成の文書として読み、どの環境でも動作する証拠にはしない。Sentryに単独の価格や提供日の発表はない。

\subsection*{生物透かし――SynthID Bio}
SynthID Bioは、タンパク質に電子透かしを埋め込み、実験室試験で機能を保ったという報告である\cite{synthidbio}。配列側はアミノ酸選択の誘導、構造側は微調整済みAlphaFold 3拡散部の重みへの埋め込みという仕組みで、合成した実物タンパク質で検証可能とされる。3つの標的で的中率・親和性・多様性が透かしなしと一致し、初の機能する透かしつき結合体の主張とされる。9月30日のNature論文が根拠で、10月1日のX投稿は関心の広がりを示すにすぎない。コードも試験管内データも重みも公開された。改ざん耐性は今後の課題とされる。バイオセキュリティ層の話題として読む。

\subsection*{指針文書と認可仕様}
残る2件は小さく添える。最前線モデルの学習向け安全事例集は、アライメントと封じ込めと監視に運用と事後対応を添えた指針文書であり、実証の報告ではない\cite{safetycases}。Workers OAuth Provider v1は、分割されたAPIと2026年7月28日のMCP認証、段階的スコープ引き上げの仕様の話題で、仕様への準拠と脅威への耐性は別の事柄である\cite{mcpauth}。索引の入り口まで読んだ範囲の話とする。

\begin{claimboundary}[資料と評価の境界]
5つの領域は独立であり、単一の安全性スコアには束ねない。論文の版はv2までに留め、v3の有無は言わない。構想の文書を動作の証拠に、指針の文書を実証に、仕様準拠を脅威への耐性に読み替えない。Xの日付は関心の広がりであり、根拠の日付ではない。
\end{claimboundary}
